\chapter{Introduction}\label{ch:intro}

\section{What \texorpdfstring{\nwrem}{nwrem} is}

\nwrem{} fits relational event models (REMs) in Stata. A relational event is one act from a sender to a receiver at a known moment: an e-mail, a phone call, a message in a chat, a handshake. The data are a list of such events with their times. A REM looks at the events one after the other and asks, at each moment, why the next event was between \emph{this} sender and \emph{this} receiver and not between any other pair of actors who could have interacted instead. Statistics of the history, such as how active the sender has been so far, whether the sender has written to this receiver before, or whether the receiver has just written to the sender, raise or lower the rate of each pair, and \nwrem{} estimates by how much. The model is that of Butts~\cite{butts2008}.

\nwrem{} is a Stata package of its own that builds on the \texttt{nwcommands} package. It works on event data declared with \texttt{nwset}'s \opt{eventtime()} option, reads the attributes of the actors from the ordinary Stata dataset, and returns its results in \texttt{e()}, so that \texttt{estimates}, \texttt{lrtest} and the other postestimation tools of Stata work. The likelihood, its gradient and its exact Hessian are evaluated by a compiled plugin; the estimates are the maximum of the likelihood (or the posterior mode) to the accuracy of the floating-point arithmetic.

\section{Relation to \texorpdfstring{\relevent}{relevent}}

The reference implementation of the model is the R package \relevent{}~\cite{relevent} by Carter T.\ Butts. \nwrem{} follows it in its model, in the names and definitions of its statistics and in its likelihoods, and its estimates were checked against it (Chapter~\ref{ch:validation}). \nwrem{} is nevertheless an \emph{independent reimplementation}, written in Stata, Mata and C, using the published description of the model and the documented behavior of \relevent{} (version~\releventver{}) as the specification.

Where the two programs differ in a way that matters to a user, this manual says so. The most important differences are these.
\begin{itemize}[nosep]
  \item The statistics are chosen by Stata options (\opt{nodsnd}, \opt{frpsndsnd}, \opt{covsnd(varlist)}), not by a character vector of effect names. The option names are those of \relevent{} in lower case. The translation table is in Appendix~\ref{app:translation}.
  \item The event list is declared with \texttt{nwset} (Chapter~\ref{ch:start}) instead of being passed as a matrix; actors that never occur in the list are added with \texttt{nwaddnodes}; the attributes of the actors are variables of the current dataset.
  \item The statistic \opt{ospsnd} is the symmetric statistic that its definition describes, not the asymmetric one that \relevent{}~\releventver{} computes (Section~\ref{sec:triadic}).
  \item \nwrem{} always computes the covariance matrix (the exact observed information); \relevent{} does so only with \texttt{hessian=TRUE}.
\end{itemize}

\section{What is covered}

\begin{itemize}[nosep]
  \item \textbf{Likelihoods.} The ordinal partial likelihood and the waiting-time (interval) likelihood, each with maximum likelihood or with the posterior mode under independent Student $t$ priors, and the conditioning on the first events (Chapter~\ref{ch:model}).
  \item \textbf{Effects.} All statistics of \relevent{} that are not identically zero: the degree and inertia effects, the recency effects, the covariate effects (of actors and of pairs, with several covariates per option), the triadic and shared-partner effects, the fixed effects of the actors and six participation shifts (Chapter~\ref{ch:effects}).
  \item \textbf{Risk sets.} The complete risk set of all ordered pairs, a restriction to given pairs, and actors that enter and leave the observation.
  \item \textbf{Postestimation.} The results of every fit are stored for \texttt{estimates} (including \texttt{estimates stats}), \texttt{lrtest} and the other standard tools (Chapter~\ref{ch:results}).
\end{itemize}

\paragraph{What \relevent{} has that \nwrem{} does not.} Bayesian simulation (MCMC and sampling importance resampling), the sampling of the risk set (\texttt{sample.size}), the simulation of event sequences, goodness-of-fit residuals and covariates that change over time are not implemented (Chapter~\ref{ch:limits}). There is no intercept, for a reason that lies in the model (Section~\ref{sec:nointercept}).

\section{Installation and requirements}\label{sec:install}

\nwrem{} needs Stata~17 or later and the core package \texttt{nwcommands}, version~3.24.0 or later. Install the core first, then \nwrem{}:
\begin{lstlisting}
. net install nwcommands, from(https://nwcommands.org/net/) replace
. nwinstall, modelpkg(nwrem)
\end{lstlisting}
\texttt{nwinstall, modelpkg(nwrem)} installs the latest release; \texttt{nwinstall, everything} installs the core and all four model packages in one step. Without \texttt{nwinstall} the package can be installed directly from its repository, and a particular version from the release tag:
\begin{lstlisting}
. net install nwrem, from(https://raw.githubusercontent.com/thomas-u-grund/nwrem/main/) replace
. net install nwrem, from(https://raw.githubusercontent.com/thomas-u-grund/nwrem/v1.5.0/) replace
\end{lstlisting}
\texttt{net describe nwrem} shows the version and the distribution date of the installed package, \texttt{nwrem version} the version of the package, of the core, the plugin in use and the version of \relevent{} that the results were compared with, and \texttt{e(version)} holds the version in the stored results. Quote it when you report an analysis.

\paragraph{The native plugin.} Every fit runs in a compiled plugin (Chapter~\ref{ch:technical}). The package ships the plugin for macOS (universal), Linux and Windows; there is no slower fallback in Mata. If the file for your platform is missing, \nwrem{} stops with return code 601 and names the file; if the installed plugin does not fit the \texttt{.ado} files or cannot be loaded, it stops with return code 3495. Reinstall the package in both cases. If the plugin does not load on your computer, you can compile it there:
\begin{lstlisting}
. nwinstall, buildplugin(nwrem)
\end{lstlisting}
This needs a C compiler (on macOS the command line tools of Xcode, on Linux \texttt{gcc} and \texttt{make}, on Windows the Microsoft Build Tools); restart Stata afterwards. The C source is part of the package (\texttt{net get nwrem, from(https://raw.githubusercontent.com/thomas-u-grund/nwrem/main)} copies it into the current directory) and of the repository \url{https://github.com/thomas-u-grund/nwrem}.

\paragraph{The Mata library.} The Mata library of \nwrem{} (\texttt{lnwrem.mlib}) is built with Stata~17, the oldest supported version, because a library built by a newer Stata is ignored by an older one (return code 3499). If the library does not load in your Stata, build it with your own Stata: \texttt{nwinstall, buildmata(nwrem)}. The command clears the data and the Mata objects in memory, so save first, and restart Stata afterwards.

\paragraph{Platforms and versions.} Stata~17 or later is required; Stata~16 and older have not been tested and are not expected to work. The full test suite of version~\pkgversion{} was run on macOS with Stata~17, Stata~18 and StataNow~19.5. The timings in Chapter~\ref{ch:technical} were measured on Linux (Ubuntu 22.04, StataNow~19.5 MP). The test suite of version~1.4.0 was also run on Linux; for version~\pkgversion{} it has not yet been run on Linux or Windows, and Windows has never been tested.

\paragraph{The examples of this manual.} All listings are the output of Stata~17 runs with \nwrem{}~\pkgversion{} and the core package~\corever{}; the do-files that produce them are in the folder \texttt{manual/ex} of the repository. The data are the e-mail log of a manufacturing company (Section~\ref{sec:data}); they are distributed with the core package and read by \texttt{nwwebuse}.

\section{License}

\nwrem{} is free software under the GNU General Public License, version~3, and Thomas U.\ Grund is its sole author. Because the package follows the structure of \relevent{}, whose license is GPL ($\geq 2$), it is distributed under the GPL-3, which is compatible with it. The example data are not covered by the GPL-3 of the code; their source and license are given in Section~\ref{sec:data}.
